\subsection{CAPS-Inspired Formulation of Trait-Conditioned C--K Role-Playing}
\label{sec:caps_formulation}

The central challenge of professional role-playing is not to reproduce an expert's background information, but to simulate the individually distinctive reasoning process through which an expert responds to conflict. Design provides a suitable setting for this problem because design tasks are rarely well-defined, closed-form problems with predetermined objectives, constraints, and evaluation criteria. Instead, they involve wicked tensions among technical feasibility, aesthetic quality, commercial interests, user well-being, and social responsibility. Even experts with comparable domain knowledge may decompose the same conflict differently, seek different forms of evidence, explore different conceptual possibilities, and draw different boundaries for compromise.

CAPS theory characterizes individual differences not only as stable average tendencies, but also as relatively stable \emph{if--then} patterns of behavioral variation under situational features \citep{mischel1995cognitive}. We use this account as a modeling motivation rather than as a causal claim about the measured traits. Specifically, we treat a designer's V30 profile as a stable conditioning variable whose association with observable C--K operations may vary with the design situation.

Let $V_i$ denote the 30-dimensional professional trait profile of designer $i$, and let $X$ denote a public design situation containing multiple conflicts. We formulate role-playing as the joint generation of an auditable C--K trajectory $M_i=(m_1,\ldots,m_T)$ and a first-person decision and justification $Y_i$:
\begin{equation}
    (V_i,X) \mapsto (M_i,Y_i).
    \label{eq:roleplay_formulation}
\end{equation}
The paired output is important: $M_i$ makes the sequence of observable reasoning operations available for inspection, whereas $Y_i$ expresses the designer's situated response. The trajectory is not a transcription of a model's hidden chain of thought.

At the level of individual operations, the conceptual conditional relationship is
\begin{equation}
    p(m_t \mid m_{<t},X,V_i),
    \label{eq:trait_conditioned_transition}
\end{equation}
which states that the profile can alter the relative propensity of different C--K operations in a given situation, rather than deterministically fixing a response across all situations. The expression describes the target conditional structure; it does not require a separate planner that first generates $M_i$ and then a second model that realizes $Y_i$. As detailed in Section~\ref{sec:trait_conditioned_cktm}, TC-CKTM instead returns both fields in a single joint generation.

This formulation separates stable dispositions, situational conditions, observable reasoning operations, and textual responses. V30 is an individual-level profile rather than a transient C--K mode, while C--K theory represents the unfolding reasoning process. Accordingly, we do not treat problem framing or trait/value conditions as additional states in the C--K trajectory. The full implementation subsequently augments $(V_i,X)$ with professional context, prior Q2 reasoning records, frozen exploratory routing rules, and controlled cross-user examples; these channels provide soft, auditable conditioning without overriding the facts of the public situation.
